\chapter{What does an economy keep alive?}\label{ch:motivation}

\section{Why people made money}
No household can make everything it needs, and no community can rely on everyone needing the same thing at the same moment. People have therefore had to exchange, remember obligations, compare unlike contributions and settle claims across time. The solutions did not appear in one place on one day. Historical records include accounts kept on clay tablets, measured commodities, weighed metal and, later, coins. Money has taken several forms; its history is not a simple story in which coins suddenly replaced a world of direct barter.\cite{intro-money-history}

Whatever its form, money can make a remarkable kind of cooperation possible. It gives people a common unit in which to state an obligation, a means of paying it and a way to carry a claim into the future. A carpenter can be paid without having to take precisely the food that the customer happens to have that day. A hospital can buy supplies from people it has never met. These functions depend on trust and institutions, but they free cooperation from the need for every exchange to be settled by an immediate swap.\cite{intro-money-functions} We should begin with that achievement, not with the claim that money was a mistake.

The question is what a monetary account tells us \emph{after} a payment succeeds. Suppose a town's clinic has paid its electricity bill and ordered medicine, but no water reaches its treatment room because a pipe has failed. The payment establishes that a claim was discharged under an agreement. The service requires a working chain of reservoirs, pipes, power, maintenance and human attention. If one link fails, a financially completed transaction can coexist with incomplete physical provision. If someone repairs the pipe without charging anyone, the reverse can occur: a vital service returns without a sale.

There is no paradox. Money records and coordinates claims; water, care and functioning equipment are conditions in the world. An economy should make it possible for people to live, recover, care for one another and hand a functioning world to those who follow them. How can it tell whether it is doing so? The word \emph{tell} does not require a central mind. It asks what information an arrangement records and what responses that information encourages. A clinic may monitor its water, a town may protect a catchment, and a government may regulate a pollutant. The ordinary monetary result of an exchange does not, by itself, report whether these life-supporting physical conditions improved or deteriorated.

\section{A measure can be excellent at the wrong question}
The versatility that makes money useful also limits what one amount can report. A sum can travel across goods, people and years precisely because it does not describe every physical property of the underlying transaction. Two products sold for the same price can require different material inputs, leave different wastes and contribute very differently to health or resilience. Equal revenue is not equal physical consequence.

An enterprise deciding whether to repair a pipe faces costs and expected returns. If it cannot recover the repair cost, the financially rational choice for that enterprise may be delay, even when the delay imposes a larger physical cost on households. The opposite possibility matters too: maintenance may be profitable and beneficial at the same time. There is no universal rule that financial gain destroys function. There is simply no identity that makes the sign of a financial return equal the sign of a physical change.

The mismatch can be corrected in particular cases. Law, public services, contracts, standards, taxes and community obligations can bring neglected effects into a decision. Their existence is important: the present economy is not a lawless machine whose only instruction is profit. But each correction must know which physical consequence it is trying to address. Without that knowledge, one can reward an activity because it is lucrative and discover its hidden costs only after water, soil, equipment or health has been damaged.

\section{Want, need and the power to pay}
In a market, a person's expressed demand depends partly on purchasing power. A buyer who can pay for a second luxury delivery may exert more influence on a seller than someone who urgently needs a first delivery of safe water but cannot pay. This is not an accusation against the seller. It is a distinction between a transaction's signal and a person's need. Money answers how much can be offered under a set of rights and prices; need asks what is required for a function or a life.

The distinction becomes sharper when the future has no present purchaser. A child not yet born cannot bid for a stable climate or an intact aquifer. A forest's contribution to rainfall or habitat may support many people without appearing on the receipt for a nearby harvest. These effects can be represented through institutions, but a price tag on one immediate sale does not automatically carry them.

It is tempting to answer by saying that every need should be assigned a price. Pricing can help in some circumstances, but it does not settle how the need is measured, who has authority to set the price, whether those affected can pay, or whether a loss is reversible. A financially visible problem is not necessarily a physically solved problem. We need to know what the world actually contains and what an action actually changes.

\begin{intuitionbox}{The question of this book}
Before asking whether an action has money, ask what it takes from the physical world, what it leaves there, and which continuing function depends on the change. A useful economy needs both a way to coordinate people and a way to remain answerable to those physical facts.
\end{intuitionbox}

\section{A world inherited rather than purchased}
Each generation receives a planet it did not manufacture. Sunlight arrives, water circulates, organisms reproduce, and human beings maintain the networks through which these processes become shelter, food, care and knowledge. None of this means that nature is a perfectly balanced machine or that every existing arrangement should be held still. Living systems change. A sustainable inheritance is not a frozen photograph; it is the continued possibility of life-supporting function through change.

That possibility can be lost while accounts still rise. It can also be preserved by work whose immediate monetary return is small. A repaired wetland, a maintained water line or a nurse's unpaid care may not fit neatly into a calculation centred on sales. They nonetheless matter to what people can do tomorrow. The question is not whether money is useful; it is whether financial success alone is an adequate compass for a civilisation that must keep its physical home viable.

We will eventually build a mathematical account of particular physical actions. First we must understand why the familiar account can miss them. The next chapter follows the incentive faced by an ordinary decision-maker: what does a price reward, what does it leave outside the deal, and why can maximising a return become different from maintaining the conditions on which the return depends?
